\scanpage{77}
% archive sheet number: 76
% L1-5
\uline{Lemme!} Soit $G$ un groupe algébrique connexe sur le corps algébriquement clos $\Omega$ de car. $p\ne0$, \deleted{soit $q$ une puissance de $p$} \inserted{``défini'' sur le corps $\unsure{\mathbb F_q|F_q}$ à $q$ éléments, soit} $\psi$ un automorphisme \inserted{du groupe algébrique} \deleted{de} $G$, considérons \deleted{l'endomorphisme du groupe $G$\unsure{,|.} défini par} \deleted{\unsure{abstrait}} l'application de $G$ dans $G$ donnée par
% L6
\[
\deleted{\unsure{\theta|\varphi}}\qquad x\longrightarrow x(\psi x^{(q)})^{-1}
\]
% L7-9
où $x^{(q)}$ est le transformé de $x$ par l'automorphisme qui correspond à $\lambda\longrightarrow\lambda^{(q)}$ dans $\Omega$ (dont le corps des inv. est $\mathbb F_q$) cette application est alors \uline{surjective}.

% L10-14
\deleted{On considère l'application $x\longrightarrow$}
Considérons la différentielle de cette application à l'élément neutre : par un calcul purement formel, on trouve que c'est l'identité (car la différentielle de $x\unsure{\text{---}|\longrightarrow}\psi x^{(q)}$ est la \deleted{\unsure{transformée}}\inserted{composée} de $d\psi(e)$ avec la différentielle de $x^{(q)}$, laquelle est nulle)
% L15-18
Donc l'application envisagée a pour image un \uline{voisinage} de $e$. C'est d'ailleurs un hom. croisé (quand $G$ opère sur lui $=$ par automorphismes intérieurs par $\psi x^{(q)}$ :
\[
\unsure{\Theta|\theta}(xy)=\Theta(x)\operatorname{int}(\psi x^{(q)})\Theta(y).
\]
% L19-21
Si $G$ est abélien, c'est que $\Theta(G)$ est donc un sous-groupe, qui est \inserted{ouvert donc} identique à $G$ puisque $G$ est connexe.

Dans le cas général, on \deleted{\unsure{a}}
